\chapter{Le village}
\label{chap:le-village}

C'est dans un village orc, loin des grandes routes empruntées par les marchands
et des villes où se pressaient les autres peuples, que commence cette histoire.

Le village occupait le sommet légèrement aplati d'une colline dominant une
vaste région couverte de forêts. Quelques sentiers étroits serpentaient entre
les arbres et permettaient de rejoindre les terres alentour, mais aucun n'était
suffisamment fréquenté pour attirer régulièrement les caravanes. De loin, on
distinguait surtout la palissade de troncs qui encerclait le sommet, quelques
toits irréguliers dépassant derrière elle et les fines colonnes de fumée qui
s'élevaient des foyers. Tout autour, la forêt descendait sur les flancs de la
colline avant de reprendre possession du paysage aussi loin que portait le
regard.

À l'intérieur de la palissade, les habitations s'étaient développées au fil des
années sans suivre de véritable plan. Elles étaient construites de bois, de
pierre, de peaux et de tout ce que les Orcs avaient pu tirer des environs ou
rapporter de leurs expéditions. Des chemins de terre circulaient entre les
maisons, les réserves, quelques enclos et les ateliers où l'on travaillait le
bois ou le métal. La palissade elle-même était robuste sans chercher à devenir
une véritable fortification : elle devait avant tout tenir les bêtes sauvages
à distance et rendre plus difficile l'approche d'un groupe hostile.

Au centre du village se trouvait son élément le plus remarquable. Le sommet de
la colline y avait été profondément excavé, jusqu'à former une vaste fosse
irrégulière dont les parois mêlaient terre compactée et roche nue. Quelques
aménagements de bois et de pierre en bordaient la partie supérieure afin de
permettre aux habitants d'observer ce qui se déroulait en contrebas, mais
l'essentiel de l'arène avait été façonné directement dans le terrain. Ses
parois abruptes étaient suffisamment hautes pour qu'un combattant, et surtout
une créature qu'on y avait enfermée, ne puisse facilement en sortir.

Cette arène occupait une place importante dans la vie du village. Les guerriers
s'y affrontaient régulièrement, autant pour s'entraîner que pour mesurer leurs
progrès, et ces combats attiraient volontiers les autres habitants. On pouvait
également y voir un Orc suffisamment téméraire affronter une créature capturée
dans les bois ou sur les terres voisines. Avec le temps, ces affrontements
étaient devenus davantage qu'un simple entraînement : ils constituaient aussi
un divertissement, un événement autour duquel une partie du village se
rassemblait.

La vie y était rude, mais elle ne l'était pas davantage que les Orcs qui y
habitaient. On y chassait, on élevait quelques animaux, on travaillait le bois,
la pierre et le métal, et chacun participait d'une manière ou d'une autre à la
survie du clan. La force y était respectée, sans pour autant se résumer à celle
des muscles : savoir suivre une piste, trouver de la nourriture, supporter le
froid, combattre, reconnaître un danger ou défendre les siens comptait tout
autant. Les enfants apprenaient tôt à se débrouiller et recevaient, en
grandissant, les enseignements qui devaient leur permettre de devenir à leur
tour des membres à part entière du village.

L'isolement de la colline ne signifiait toutefois pas que personne ne
s'aventurait jamais sur les terres du clan. Des chasseurs, des voyageurs ou de
petits convois empruntaient parfois les sentiers voisins sans savoir exactement
où ils les conduiraient. Certains étaient simplement repoussés ou dépouillés ;
d'autres avaient moins de chance et étaient ramenés au village.

Les prisonniers pouvaient être contraints de travailler quelque temps pour le
clan. Certains finissaient également au fond de l'arène, opposés à une bête
capturée dans les environs ou à l'un des guerriers du village. Pour les Orcs
qui assistaient à ces affrontements, il n'y avait là rien de particulièrement
cruel ou exceptionnel. Ces pratiques existaient depuis aussi longtemps que la
plupart d'entre eux pouvaient s'en souvenir et faisaient simplement partie de
leur manière de vivre. Rares étaient ceux qui songeaient à se demander s'il
pouvait en être autrement.

\illustration
  {0100.png}
  {Le village de Khaer Ghal}
  {village}

\scenebreak

C'est dans ce village que, par une nuit comme tant d'autres, vint au monde
Khaer Ghal.

Il était le fils du chef.

Son père était un Orc imposant, respecté autant pour sa force que pour
l'autorité avec laquelle il dirigeait les siens. Il avait gagné sa place depuis
longtemps et voyait naturellement dans la naissance de son fils la continuité
de sa famille au sein du clan. Lorsque Khaer Ghal serait suffisamment grand,
il apprendrait à chasser, à se battre, à survivre et à défendre le village
comme les autres enfants. Davantage encore, son père attendrait de lui qu'il se
montre digne du nom et de la place dont il héritait.

Sa mère regardait les choses différemment. Non qu'elle rejetât leur peuple, ses
coutumes ou la vie qu'ils menaient sur cette colline ; elle en faisait
pleinement partie. Mais là où son père songeait déjà à ce que leur fils devrait
un jour devenir, elle semblait davantage attentive à l'enfant qu'elle tenait
dans ses bras. Elle lui transmettrait elle aussi ce qu'elle savait et ne
chercherait jamais à le soustraire à l'éducation des siens. Certains de ses
gestes, certaines de ses paroles ou simplement certains de ses silences
laissaient pourtant parfois entrevoir une autre manière de regarder le monde.

Pour l'heure, toutes ces attentes appartenaient encore à un avenir lointain.
Khaer Ghal n'était qu'un nouveau-né, minuscule lorsqu'on le comparait aux
larges bras orcs qui le portaient. Sa peau possédait déjà cette teinte
gris-bleu qui resterait la sienne en grandissant et, lorsque ses paupières
s'ouvraient suffisamment longtemps, elles révélaient des yeux d'un jaune doré
étonnamment vif. Ses oreilles, déjà légèrement pointues, encadraient un visage
encore rond sur lequel rien ne permettait de deviner l'Orc qu'il deviendrait.

Il ne connaissait alors du monde que la chaleur des bras qui le tenaient, les
voix graves qui l'entouraient et les bruits du village au-delà des murs : le
crépitement des foyers, le choc lointain du métal, les animaux dans leurs
enclos et le murmure du vent contre la palissade. Pour ceux qui se penchaient
sur lui cette nuit-là, il était avant tout le fils du chef, un enfant qui
grandirait parmi les siens, apprendrait leurs coutumes et trouverait un jour sa
place dans le clan.

Ce jour était encore loin.

Pour l'heure, Khaer Ghal venait simplement de naître.